\documentclass[11pt]{article}
\usepackage[margin=0.65in]{geometry}
\usepackage{booktabs}
\usepackage{float}
\usepackage[T1]{fontenc}
\usepackage[utf8]{inputenc}
\begin{document}
\section*{Input-output price shocks and inflation inequality}
\begin{table}[H]
\centering
\caption{Cumulative Energy Shock Exposure, 2020--2024}
\scriptsize
\setlength{\tabcolsep}{4.2pt}
\begin{tabular}{lrrrrr}
\toprule
Country & \shortstack{Average inflation\\Total} & \shortstack{Average inflation\\Direct} & \shortstack{Income inflation\\gap} & \shortstack{Direct income\\inflation gap} & \shortstack{Indirect income\\inflation gap}
 \\
\midrule
DE & 11.1 & 4.3 & 1.4 & 1.8 & -0.4
 \\
FR & 16.5 & 4.7 & 1.0 & 1.6 & -0.5
 \\
IT & 11.9 & 4.0 & -0.9 & -0.6 & -0.3
 \\
ES & 3.2 & 1.4 & 0.7 & 0.6 & 0.1
 \\
NL & 9.0 & 3.6 & -0.3 & 0.6 & -0.9
 \\
BE & 16.8 & 7.6 & 0.8 & 2.2 & -1.4
 \\
EA20 & 11.6 & 4.1 & 0.6 & 1.1 & -0.4
 \\
\bottomrule
\end{tabular}
\par\vspace{0.35em}\footnotesize\noindent Sources: Eurostat HICP, HBS, FIGARO, authors' calculations.
\par\vspace{0.25em}\footnotesize\noindent Notes. Entries are cumulative percentage-point effects of observed energy-price increases between January 2020 and December 2024. The direct shocks are country-specific cumulative HICP changes mapped to FIGARO energy sectors: upstream oil and gas extraction (B), refined petroleum products (C19), and electricity/gas utilities (D35). The price response is computed as $\Delta p = (I - A^{\prime})^{-1}s$, with $A$ built from the latest FIGARO matrix used in the simulations. The direct effect is the component from the shocked energy sectors in the household consumption basket. Indirect is the Leontief propagation component. Average inflation total includes direct and indirect effects before rounding. Income inflation gap = Q1 - Q5.
\end{table}
\vspace{0.6em}
\begin{table}[H]
\centering
\caption{Effect of a 10 Percent Shock to Cereal Prices, 2023}
\scriptsize
\setlength{\tabcolsep}{4.2pt}
\begin{tabular}{lrrrrr}
\toprule
Country & \shortstack{Average inflation\\Total} & \shortstack{Average inflation\\Direct} & \shortstack{Income inflation\\gap} & \shortstack{Direct income\\inflation gap} & \shortstack{Indirect income\\inflation gap}
 \\
\midrule
DE & 0.7 & 0.2 & -0.0 & 0.0 & -0.0
 \\
FR & 1.2 & 0.3 & -0.1 & 0.0 & -0.1
 \\
IT & 1.2 & 0.5 & 0.2 & 0.1 & 0.1
 \\
ES & 1.2 & 0.3 & -0.0 & 0.0 & -0.0
 \\
NL & 1.1 & 0.3 & -0.1 & 0.0 & -0.1
 \\
BE & 1.1 & 0.3 & -0.1 & 0.0 & -0.1
 \\
EA20 & 1.0 & 0.3 & -0.0 & 0.0 & -0.0
 \\
\bottomrule
\end{tabular}
\par\vspace{0.35em}\footnotesize\noindent Sources: Eurostat HICP, HBS, FIGARO, authors' calculations.
\par\vspace{0.25em}\footnotesize\noindent Notes. Entries are percentage-point effects of the sectoral input price shock shown in the table title. The price response is computed as $\Delta p = (I - A^{\prime})^{-1}s$, with $A$ built from FIGARO. The direct effect is the component from the shocked sector in the household consumption basket. Indirect is the Leontief propagation component. Average inflation total includes direct and indirect effects before rounding. Income inflation gap = Q1 - Q5 of revenues.
\end{table}
\vspace{0.6em}
\begin{table}[H]
\centering
\caption{Effect of a 10 Percent Shock to Electronic Component Prices, 2023}
\scriptsize
\setlength{\tabcolsep}{4.2pt}
\begin{tabular}{lrrrrr}
\toprule
Country & \shortstack{Average inflation\\Total} & \shortstack{Average inflation\\Direct} & \shortstack{Income inflation\\gap} & \shortstack{Direct income\\inflation gap} & \shortstack{Indirect income\\inflation gap}
 \\
\midrule
DE & 0.2 & 0.1 & -0.0 & -0.0 & -0.0
 \\
FR & 0.2 & 0.1 & -0.0 & -0.0 & -0.0
 \\
IT & 0.2 & 0.1 & -0.0 & -0.0 & -0.0
 \\
ES & 0.1 & 0.0 & -0.0 & -0.0 & -0.0
 \\
NL & 0.2 & 0.1 & -0.0 & -0.0 & -0.0
 \\
BE & 0.2 & 0.0 & -0.0 & -0.0 & -0.0
 \\
EA20 & 0.2 & 0.1 & -0.0 & -0.0 & -0.0
 \\
\bottomrule
\end{tabular}
\par\vspace{0.35em}\footnotesize\noindent Sources: Eurostat HICP, HBS, FIGARO, authors' calculations.
\par\vspace{0.25em}\footnotesize\noindent Notes. Entries are percentage-point effects of the sectoral input price shock shown in the table title. The price response is computed as $\Delta p = (I - A^{\prime})^{-1}s$, with $A$ built from FIGARO. The direct effect is the component from the shocked sector in the household consumption basket. Indirect is the Leontief propagation component. Average inflation total includes direct and indirect effects before rounding. Income inflation gap = Q1 - Q5 of revenues.
\end{table}
\vspace{0.6em}
\end{document}
